\documentclass[a4paper]{article}
\usepackage[utf8]{inputenc}
\usepackage[T5]{fontenc}
\usepackage[vietnamese]{babel}
\usepackage[fontsize=13pt]{scrextend}
\usepackage{mathptmx}
\usepackage{amsmath}
\usepackage{geometry}
\usepackage{setspace}
\usepackage{indentfirst}
\usepackage{graphicx}
\usepackage{xcolor}
\usepackage{booktabs}
\usepackage{longtable}
\usepackage{tabularx}
\usepackage{array}
\usepackage{enumitem}
\usepackage{listings}
\usepackage{caption}
\usepackage{float}
\usepackage{fancyhdr}
\usepackage{hyperref}
\usepackage{microtype}
\usepackage{tikz}
\usetikzlibrary{arrows.meta,positioning,fit,shapes.geometric,calc,backgrounds}

\geometry{left=3cm,right=2cm,top=2cm,bottom=2cm}
\setstretch{1.3}
\setlength{\parindent}{1cm}
\setlength{\parskip}{3pt}
\setlength{\emergencystretch}{3em}
\setlist{nosep}
\renewcommand{\arraystretch}{1.25}

\definecolor{bluec}{RGB}{24,91,166}
\definecolor{cyanc}{RGB}{0,135,165}
\definecolor{greenc}{RGB}{32,130,82}
\definecolor{orangec}{RGB}{210,116,24}
\definecolor{redc}{RGB}{180,55,55}
\definecolor{violetc}{RGB}{105,75,170}
\definecolor{grayc}{RGB}{92,100,112}
\definecolor{codebg}{RGB}{247,248,250}
\definecolor{codeframe}{RGB}{210,214,220}

\hypersetup{
  colorlinks=true,linkcolor=bluec,urlcolor=bluec,citecolor=bluec,
  pdftitle={Báo cáo mã nguồn MeshGateWay},
  pdfauthor={Nhóm phát triển MSMS}
}

\lstdefinestyle{textstyle}{
  backgroundcolor=\color{codebg},frame=single,rulecolor=\color{codeframe},
  basicstyle=\ttfamily\small,breaklines=true,columns=fullflexible,
  showstringspaces=false
}

\pagestyle{fancy}
\fancyhf{}
\fancyhead[L]{Báo cáo mã nguồn MeshGateWay}
\fancyhead[R]{ESP-IDF}
\fancyfoot[C]{\thepage}
\setlength{\headheight}{16pt}

\newcommand{\pathcode}[1]{\nolinkurl{#1}}
\newcommand{\codeword}[1]{\texttt{\detokenize{#1}}}
\tikzset{
  block/.style={draw=bluec,rounded corners=3pt,thick,align=center,
    minimum height=9mm,fill=blue!5,font=\small},
  task/.style={draw=violetc,rounded corners=3pt,thick,align=center,
    minimum height=9mm,fill=violet!8,font=\small},
  store/.style={draw=orangec,thick,align=center,fill=orange!9,
    cylinder,shape border rotate=90,aspect=.25,font=\small},
  decision/.style={draw=orangec,thick,diamond,aspect=2,align=center,
    fill=orange!8,font=\small},
  flow/.style={-{Latex[length=2.4mm]},thick,draw=bluec},
  dataflow/.style={-{Latex[length=2.4mm]},thick,draw=greenc},
  warnflow/.style={-{Latex[length=2.4mm]},thick,dashed,draw=redc},
  note/.style={draw=grayc,rounded corners=2pt,fill=gray!8,
    align=left,font=\scriptsize}
}

\begin{document}
\pagenumbering{Alph}
\begin{titlepage}
  \centering
  \vspace*{1.2cm}
  {\large BÁO CÁO PHÂN TÍCH MÃ NGUỒN\par}
  \vspace{1.5cm}
  {\Huge\bfseries MeshGateWay\par}
  \vspace{0.8cm}
  {\Large Gateway UART -- Wi-Fi -- WebSocket của hệ thống MSMS\par}
  \vspace{2cm}
  \begin{tabular}{rl}
    Nền tảng: & ESP32 / ESP-IDF / FreeRTOS \\
    Giao tiếp: & UART, Wi-Fi, HTTP, mDNS, WebSocket \\
    Giao diện: & LVGL, captive portal \\
    Phiên bản tài liệu: & 1.0 \\
  \end{tabular}
  \vfill
  {\large Tài liệu được xây dựng từ trạng thái mã nguồn hiện tại\par}
  \vspace{0.3cm}
  {\large Tháng 07 năm 2026\par}
\end{titlepage}

\pagenumbering{roman}
\tableofcontents
\clearpage
\listoffigures
\clearpage
\listoftables
\clearpage
\pagenumbering{arabic}

\section{Mục đích và phạm vi}

\codeword{MeshGateWay} là firmware trung gian giữa nút root của mạng cảm biến và
hệ thống máy chủ. Gateway không trực tiếp vận hành mạng mesh trong luồng chính
hiện tại. Nó nhận chuỗi telemetry từ root qua UART, bổ sung trạng thái của chính
gateway, rồi chuyển dữ liệu lên backend bằng WebSocket.

Tài liệu tập trung vào mã đang được gọi từ \codeword{app_main()}. Những module có
trong cây mã nhưng chưa được nối vào luồng chạy, chẳng hạn FOTA hoàn chỉnh,
\codeword{ESP_Bus} và một số tiện ích cũ, được ghi rõ là hạ tầng dự phòng thay vì
mô tả như chức năng đã hoạt động.

\section{Bối cảnh và ranh giới hệ thống}

\subsection{Sơ đồ nguồn và đích dữ liệu}

Hình~\ref{fig:context} trả lời ba câu hỏi: dữ liệu sinh ra ở đâu, gateway xử lý
gì, và dữ liệu được tiêu thụ ở đâu. Mũi tên xanh lá là luồng telemetry; mũi tên
xanh dương là điều khiển hoặc cấu hình.

\begin{figure}[H]
\centering
\resizebox{.96\textwidth}{!}{
\begin{tikzpicture}[node distance=12mm and 18mm]
  \node[block,minimum width=3.2cm,minimum height=1.5cm,fill=green!8] (sensors)
    {\textbf{Các node cảm biến}\\BME/DHT/PM/CO$_2$\\JSON có trường \codeword{M}};
  \node[block,right=of sensors,minimum width=3.2cm,minimum height=1.5cm] (root)
    {\textbf{MSMS root}\\nhận dữ liệu mesh\\đẩy dòng qua UART};
  \node[block,right=of root,minimum width=4.2cm,minimum height=1.8cm,fill=cyan!8] (gw)
    {\textbf{MeshGateWay ESP32}\\ghép dòng và lọc UART\\tạo \codeword{uart_rx} và
    \codeword{gateway_status}};
  \node[block,right=of gw,minimum width=3.4cm,minimum height=1.5cm,fill=orange!8] (server)
    {\textbf{Backend}\\WebSocket cổng 9090\\phân loại và lưu dữ liệu};
  \node[block,below=16mm of server,minimum width=3.4cm,fill=violet!8] (web)
    {\textbf{Dashboard}\\hiển thị gateway, node\\và dữ liệu cảm biến};
  \node[block,below=16mm of gw,minimum width=4.2cm,fill=yellow!12] (config)
    {\textbf{Người vận hành}\\captive portal / HTTP\\chọn Wi-Fi và URL WS};

  \draw[dataflow] (sensors) -- node[above,font=\scriptsize]{mesh / UDP} (root);
  \draw[dataflow] (root) -- node[above,align=center,font=\scriptsize]{UART\\115200} (gw);
  \draw[dataflow] (gw) -- node[above,align=center,font=\scriptsize]{local: \codeword{ws://}\\server: \codeword{wss://}} (server);
  \draw[dataflow] (server) -- node[right,font=\scriptsize]{API / WS} (web);
  \draw[flow] (config) -- node[right,font=\scriptsize]{HTTP :80} (gw);
  \draw[warnflow,bend left=18] (server) to node[below,align=center,font=\scriptsize]{lệnh / version\\(xử lý sơ bộ)} (gw);
\end{tikzpicture}}
\caption{Bối cảnh và đường đi đầu cuối của dữ liệu}
\label{fig:context}
\end{figure}

Gateway có hai nguồn dữ liệu độc lập:
\begin{itemize}
  \item \textbf{Nguồn từ node:} chuỗi JSON đi qua UART, được giữ nguyên trong
  trường \codeword{payload} của bản tin \codeword{uart_rx}.
  \item \textbf{Nguồn nội tại gateway:} CPU, RAM, pin, uptime, SSID, RSSI, IP và
  gateway mạng được đo tại gateway và gửi bằng \codeword{gateway_status}.
\end{itemize}

\subsection{Sơ đồ khối phần cứng và giao tiếp}

\begin{figure}[H]
\centering
\resizebox{.96\textwidth}{!}{
\begin{tikzpicture}[node distance=8mm and 14mm]
  \node[block,minimum width=4cm,minimum height=5.2cm,fill=blue!6] (esp) {};
  \node[font=\bfseries,anchor=north] at ([yshift=-3mm]esp.north) {ESP32 MeshGateWay};
  \node[block,minimum width=2.8cm,above left=3mm and 18mm of esp] (root) {MSMS root\\UART TX/RX};
  \node[block,minimum width=2.8cm,left=18mm of esp] (battery) {Pin / nguồn\\ADC1 CH6, GPIO34};
  \node[block,minimum width=2.8cm,below left=3mm and 18mm of esp] (rtc) {RTC DS3231\\I$^2$C};
  \node[block,minimum width=2.8cm,above right=3mm and 18mm of esp] (wifi) {Access Point\\Wi-Fi STA/AP};
  \node[block,minimum width=2.8cm,right=18mm of esp] (display) {Màn hình\\LVGL};
  \node[block,minimum width=2.8cm,below right=3mm and 18mm of esp] (sd) {Thẻ SD\\dung lượng};

  \node[block,minimum width=2.8cm] at ([yshift=10mm]esp.center) (uart) {UART1\\TX17, RX16};
  \node[block,minimum width=2.8cm,below=4mm of uart] (runtime) {FreeRTOS\\DataManager};
  \node[block,minimum width=2.8cm,below=4mm of runtime] (net) {TCP/IP\\HTTP + WS};

  \draw[flow,<->] (root) -- node[above,font=\scriptsize]{115200, 8N1} (esp.west |- root);
  \draw[flow] (battery) -- (esp.west |- battery);
  \draw[flow,<->] (rtc) -- (esp.west |- rtc);
  \draw[flow,<->] (esp.east |- wifi) -- (wifi);
  \draw[flow] (esp.east |- display) -- (display);
  \draw[flow,<->] (esp.east |- sd) -- (sd);
\end{tikzpicture}}
\caption{Sơ đồ khối phần cứng quanh gateway}
\end{figure}

\section{Tổ chức mã nguồn}

\begin{longtable}{p{5.3cm}p{8.7cm}}
\caption{Các module quan trọng của MeshGateWay}\\
\toprule
\textbf{Đường dẫn} & \textbf{Trách nhiệm} \\
\midrule
\endfirsthead
\toprule
\textbf{Đường dẫn} & \textbf{Trách nhiệm} \\
\midrule
\endhead
\pathcode{main/MeshGateWay.c} & Điểm vào, khởi tạo NVS, SPIFFS, phần cứng và các task.\\
\pathcode{components/core/Datamanager} & Các cấu trúc trạng thái dùng chung cho phần cứng, metrics, Wi-Fi, WS, UART và telemetry.\\
\pathcode{components/core/UartToNode} & Bắt tay root, đọc UART, lọc và đẩy dữ liệu vào queue.\\
\pathcode{components/core/WebSocKetHandle} & Chọn URL, phân giải mDNS, duy trì client WS/WSS, gửi dữ liệu.\\
\pathcode{components/core/WifiManager} & STA/AP, captive portal, DNS, HTTP và mDNS \codeword{meshgateway.local}.\\
\pathcode{components/core/SystemMonitor} & Thu thập metrics mỗi giây.\\
\pathcode{components/core/ScreenManager} & LVGL, hiển thị trạng thái và biểu đồ RX/TX.\\
\pathcode{components/core/PowerManager} & Đọc ADC và ước lượng điện áp/phần trăm pin.\\
\pathcode{components/drivers/SD_Card} & Khởi tạo và đọc dung lượng thẻ SD.\\
\pathcode{components/core/FOTAManager} & Mã OTA cũ/chưa được nối vào luồng chạy hiện tại.\\
\pathcode{components/core/WebConfigWifi} & Tệp HTML được đóng gói vào SPIFFS cho captive portal.\\
\bottomrule
\end{longtable}

\section{Kiến trúc phần mềm nội bộ}

\subsection{Quan hệ giữa các task và trạng thái dùng chung}

\begin{figure}[H]
\centering
\resizebox{.94\textwidth}{!}{
\begin{tikzpicture}[node distance=17mm and 18mm]
  \node[store,minimum width=2.8cm] (queue) {Queue UART\\8 $\times$ 512 B};
  \node[task,below left=of queue,minimum width=3.4cm] (uart) {\textbf{uart\_to\_node}\\đọc/ghi UART};
  \node[task,below right=of queue,minimum width=3.4cm] (ws) {\textbf{ws\_hdl}\\WebSocket};
  \node[task,right=30mm of ws,minimum width=3.4cm] (wifi) {\textbf{wifi\_connect}\\\textbf{wifi\_mgr}};

  \node[block,below=19mm of ws,minimum width=8.2cm,minimum height=1.9cm,fill=cyan!7] (dm)
    {\textbf{DataManager -- các cấu trúc toàn cục}\\
    \codeword{g_metrics}, \codeword{g_telemetry}, \codeword{g_wifi},
    \codeword{g_ws}, \codeword{g_uart}, \codeword{g_hw}};
  \node[task,below left=18mm and 5mm of dm,minimum width=3.4cm] (monitor)
    {\textbf{system\_monitor}\\ghi metrics mỗi 1 s};
  \node[task,below right=18mm and 5mm of dm,minimum width=3.4cm] (screen)
    {\textbf{lvgl\_task}\\đọc snapshot, UI 250 ms};
  \node[store,right=24mm of dm,minimum width=3.1cm,minimum height=1.2cm] (spiffs)
    {NVS\\SPIFFS};

  \draw[dataflow] (uart) -- node[left,font=\scriptsize]{ghi item} (queue);
  \draw[dataflow] (queue) -- node[right,font=\scriptsize]{đọc item} (ws);
  \draw[flow,<->] (uart) -- (dm);
  \draw[flow,<->] (ws) -- (dm);
  \draw[flow,<->] (wifi) -- (dm);
  \draw[dataflow] (monitor) -- node[left,font=\scriptsize]{mutex} (dm);
  \draw[dataflow] (dm) -- node[right,font=\scriptsize]{snapshot} (screen);
  \draw[flow,<->] (wifi) -- (spiffs);
\end{tikzpicture}}
\caption{Task, queue và vùng trạng thái dùng chung}
\label{fig:tasks}
\end{figure}

\codeword{g_metrics} được bảo vệ bằng mutex. Các biến đếm telemetry được nhiều
task truy cập trực tiếp; nếu sau này tăng mức đồng thời hoặc yêu cầu số liệu tuyệt
đối, nên chuyển các bộ đếm sang thao tác atomic hay mutex riêng.

\subsection{Trình tự khởi động}

\begin{figure}[p]
\centering
\resizebox{.88\textwidth}{!}{
\begin{tikzpicture}[node distance=5mm]
  \node[block,minimum width=8.5cm] (start) {\codeword{app_main()} bật nguồn màn hình};
  \node[block,below=of start,minimum width=8.5cm] (nvs) {Khởi tạo NVS; xóa và tạo lại nếu sai phiên bản/hết trang};
  \node[block,below=of nvs,minimum width=8.5cm] (spiffs) {Mount SPIFFS tại \codeword{/spiffs}; không tự format khi lỗi};
  \node[block,below=of spiffs,minimum width=8.5cm] (state) {Gắn trạng thái Wi-Fi/WS; tạo queue 8 item và mutex metrics};
  \node[block,below=of state,minimum width=8.5cm] (hw) {Khởi động LVGL, SD, RTC DS3231 và ADC pin};
  \node[block,below=of hw,minimum width=8.5cm] (mon) {Tạo \codeword{system_monitor} -- priority 5, core 0};
  \node[block,below=of mon,minimum width=8.5cm] (wf) {\codeword{wifi_init_sta()}: thử Wi-Fi, AP/captive portal, task nền};
  \node[block,below=of wf,minimum width=8.5cm] (wst) {Tạo \codeword{ws_hdl}; stack yêu cầu từ app: 4096 B};
  \node[block,below=of wst,minimum width=8.5cm] (ut) {Gắn queue/trạng thái và khởi động task UART};
  \node[block,below=of ut,minimum width=8.5cm] (sntp) {SNTP \codeword{pool.ntp.org}, múi giờ ICT-7; chờ tối đa khoảng 28 s};
  \node[block,below=of sntp,minimum width=8.5cm] (delay) {Chờ thêm 10 s để tách tải bắt tay mạng};
  \node[block,below=of delay,minimum width=8.5cm] (ins) {Khởi tạo ESP Insights; sau đó xóa task \codeword{app_main}};
  \foreach \a/\b in {start/nvs,nvs/spiffs,spiffs/state,state/hw,hw/mon,mon/wf,wf/wst,wst/ut,ut/sntp,sntp/delay,delay/ins}
    \draw[flow] (\a) -- (\b);
\end{tikzpicture}}
\caption{Lưu đồ khởi động thực tế}
\label{fig:boot}
\end{figure}

Điểm đáng chú ý là \codeword{wifi_init_sta()} thực hiện thử kết nối đồng bộ trước
khi trả về. Vì vậy thời gian boot có thể tăng đáng kể khi nhiều SSID đều không
khả dụng. SNTP cũng chờ trong \codeword{app_main()}, nhưng các task monitor, Wi-Fi,
WebSocket và UART đã được tạo trước đó nên vẫn có thể chạy.

\section{Luồng UART giữa root và gateway}

\subsection{Cấu hình và giao thức bắt tay}

UART mặc định dùng UART1, 115200 baud, 8N1, TX GPIO17, RX GPIO16, không dùng
RTS/CTS. Driver có RX buffer 2048 B, TX buffer 1024 B và event queue 64 phần tử.

\begin{figure}[H]
\centering
\resizebox{.94\textwidth}{!}{
\begin{tikzpicture}[node distance=22mm]
  \node[block,minimum width=4.8cm,minimum height=1.5cm,fill=orange!8] (start)
    {\textbf{START\_ROOT}\\mỗi 500 ms gửi\\\texttt{Start Root} + CRLF};
  \node[block,right=of start,minimum width=4.8cm,minimum height=1.5cm,fill=green!8] (conn)
    {\textbf{CONNECTED}\\mỗi 500 ms gửi\\\texttt{Connected} + CRLF};
  \draw[flow,bend left=18] (start) to node[above,align=center,font=\scriptsize]
    {rolling buffer 48 B thấy đúng chuỗi\\\codeword{Switching to root mode...}} (conn);
  \draw[warnflow,bend left=18] (conn) to node[below,align=center,font=\scriptsize]
    {không nhận byte trong 2000 ms\\reset quy trình bắt tay} (start);
  \draw[flow,loop left] (start) to node[left,font=\scriptsize]{chưa đủ mẫu} (start);
  \draw[flow,loop right] (conn) to node[right,font=\scriptsize]{có dữ liệu UART} (conn);
\end{tikzpicture}}
\caption{Máy trạng thái bắt tay UART}
\label{fig:uartstate}
\end{figure}

Chuỗi xác nhận phải khớp chính xác
\texttt{Switching to root mode...} và kết thúc bằng CRLF. Rolling
buffer cho phép nhận chuỗi qua nhiều event UART. Sau khi kết nối, nếu đường UART
im lặng quá 2 giây, gateway quay lại \codeword{START_ROOT}.

\subsection{Đường ống xử lý dữ liệu UART}

\begin{figure}[p]
\centering
\resizebox{\textwidth}{!}{
\begin{tikzpicture}[node distance=8mm and 8mm]
  \node[block,minimum width=2.5cm] (wire) {Root UART\\dòng JSON};
  \node[block,right=of wire,minimum width=2.8cm] (event) {UART event\\RX buffer 2048 B};
  \node[decision,right=of event] (filter) {\codeword{No Node}?};
  \node[block,above=10mm of filter,minimum width=2.4cm,fill=red!7] (drop1) {Bỏ qua\\không uplink};
  \node[decision,right=of filter] (online) {WS đã\\connected?};
  \node[block,above=10mm of online,minimum width=2.4cm,fill=red!7] (drop2) {Không lưu offline\\bỏ dữ liệu};
  \node[block,right=of online,minimum width=2.8cm] (chunks) {Cắt item\\tối đa 512 B};
  \node[store,right=of chunks,minimum width=2.6cm] (queue) {Queue\\8 item};

  \node[block,below=17mm of queue,minimum width=3.0cm] (join) {Ghép đến ký tự LF\\buffer dòng 2048 B};
  \node[block,left=of join,minimum width=3.0cm] (wrap) {Escape JSON\\bọc \codeword{uart_rx}};
  \node[block,left=of wrap,minimum width=2.7cm,fill=green!8] (send) {Gửi WebSocket\\và tăng TX};

  \draw[dataflow] (wire) -- (event);
  \draw[dataflow] (event) -- (filter);
  \draw[warnflow] (filter) -- node[right,font=\scriptsize]{có} (drop1);
  \draw[dataflow] (filter) -- node[above,font=\scriptsize]{không} (online);
  \draw[warnflow] (online) -- node[right,font=\scriptsize]{không} (drop2);
  \draw[dataflow] (online) -- node[above,font=\scriptsize]{có} (chunks);
  \draw[dataflow] (chunks) -- (queue);
  \draw[dataflow] (queue) -- (join);
  \draw[dataflow] (join) -- (wrap);
  \draw[dataflow] (wrap) -- (send);
\end{tikzpicture}}
\caption{Đường ống từ byte UART tới bản tin WebSocket}
\label{fig:uartpipe}
\end{figure}

Queue có sức chứa tải danh nghĩa tối đa $8 \times 512 = 4096$ byte, chưa tính
metadata. Queue này chỉ giảm chênh lệch tốc độ ngắn hạn; nó không phải kho dữ
liệu offline. Nếu WebSocket chưa kết nối, hoặc queue đầy, telemetry có thể mất.

\begin{lstlisting}[style=textstyle,caption={Dạng bản tin uplink UART}]
{
  "type": "uart_rx",
  "len": 127,
  "payload": "{\"v\":1,\"seq\":88,\"n\":2,\"M\":\"B4:BF:E9:34:0C:CC\",...}"
}
\end{lstlisting}

Trường \codeword{M} trong payload là định danh MAC của node gửi. Gateway không
đổi \codeword{M}; backend cần parse lớp JSON bên ngoài, sau đó parse tiếp chuỗi
\codeword{payload} để phân tách node và các điểm đo.

\section{Wi-Fi, captive portal và mDNS}

\subsection{Máy trạng thái kết nối Wi-Fi}

\begin{figure}[H]
\centering
\resizebox{.82\textwidth}{!}{
\begin{tikzpicture}[node distance=12mm and 24mm]
  \node[block,minimum width=3.4cm] (boot) {Boot\\nạp cấu hình NVS};
  \node[block,right=of boot,minimum width=4.0cm] (try) {Thử từng SSID\\tối đa 5 lần/SSID\\chờ tối đa 20 s/lần};
  \node[decision,right=of try] (ok) {Có IP?};
  \node[block,above right=16mm and 27mm of ok,minimum width=3.8cm,fill=green!8] (sta)
    {STA connected\\mDNS hoạt động};
  \node[block,below right=16mm and 27mm of ok,minimum width=3.8cm,fill=orange!8] (ap)
    {AP cấu hình\\192.168.4.1\\DNS + HTTP};
  \node[block,right=25mm of ap,minimum width=4.0cm] (form)
    {Người dùng gửi\\SSID/password/WS URL};

  \draw[flow] (boot) -- (try);
  \draw[flow] (try) -- (ok);
  \draw[dataflow] (ok) -- node[above left,font=\scriptsize]{có} (sta);
  \draw[warnflow] (ok) -- node[below left,font=\scriptsize]{không} (ap);
  \draw[flow] (ap) -- (form);
  \draw[flow] (form) |- node[pos=.25,right,font=\scriptsize]{yêu cầu kết nối mới} (try.south);
  \draw[warnflow] (sta.west) -| node[pos=.25,above,font=\scriptsize]{mất kết nối} (ap.west);
  \draw[dataflow] (ap.east) -| node[pos=.25,right,align=left,font=\scriptsize]
    {STA có IP:\\tắt AP và webserver} (sta.east);
\end{tikzpicture}}
\caption{Luồng STA/AP và captive portal}
\label{fig:wifi}
\end{figure}

Khi STA lấy được IP, firmware khởi tạo mDNS với hostname
\codeword{meshgateway.local} và công bố dịch vụ HTTP cổng 80. Đây là tên của
\textbf{gateway và trang cấu hình}, khác với \codeword{systemmsems.local} là tên
máy chạy backend local.

Danh sách Wi-Fi boot hiện chứa cả credential được viết trực tiếp trong
\pathcode{wifi_manager_core.c}. Đây là rủi ro bảo mật; credential nên chuyển sang
NVS provisioning hoặc cấu hình bí mật ngoài repository.

\section{WebSocket local và server}

\subsection{Quy tắc chọn giao thức}

\begin{figure}[H]
\centering
\resizebox{.94\textwidth}{!}{
\begin{tikzpicture}[node distance=7mm]
  \node[decision] (mode) {URL bắt đầu\\bằng gì?};
  \node[block,below left=13mm and 20mm of mode,minimum width=5.1cm,fill=green!8] (local)
    {\textbf{\codeword{ws://systemmsems.local:9090}}\\mDNS tìm IPv4 trong 3 s\\
    đổi thành \codeword{ws://<IP>:9090}\\không gắn CA bundle};
  \node[block,below right=13mm and 20mm of mode,minimum width=5.1cm,fill=blue!8] (server)
    {\textbf{\codeword{wss://<server>/ws}}\\DNS thông thường\\TLS bật\\
    gắn \codeword{esp_crt_bundle_attach}};
  \node[block,below=18mm of local,minimum width=5.1cm] (plain) {WebSocket thường\\phù hợp LAN local};
  \node[block,below=18mm of server,minimum width=5.1cm] (tls) {WebSocket bảo mật\\phù hợp Internet};
  \draw[dataflow] (mode) -- node[above left,font=\scriptsize]{\codeword{ws://}} (local);
  \draw[flow] (mode) -- node[above right,font=\scriptsize]{\codeword{wss://}} (server);
  \draw[dataflow] (local) -- (plain);
  \draw[flow] (server) -- (tls);
\end{tikzpicture}}
\caption{Phân nhánh local WS và server WSS}
\label{fig:wsprotocol}
\end{figure}

Frontend chạy tại \codeword{localhost:3000} không buộc ESP32 dùng WSS. Giao thức
của ESP32 phụ thuộc URL backend được cấu hình: local hiện dùng
\codeword{ws://systemmsems.local:9090}; server Internet dùng \codeword{wss://}.
HTTP của frontend và WebSocket của gateway là hai kết nối độc lập.

\subsection{Vòng đời WebSocket}

\begin{figure}[H]
\centering
\resizebox{\textwidth}{!}{
\begin{tikzpicture}[node distance=8mm and 10mm]
  \node[block,minimum width=2.7cm] (wait) {Chờ Wi-Fi};
  \node[decision,right=of wait] (wifi) {STA có IP?};
  \node[block,right=of wifi,minimum width=3.2cm] (resolve) {Lấy URL\\resolve mDNS nếu local};
  \node[block,right=of resolve,minimum width=3.0cm] (start) {Init và start\\WS client};
  \node[block,below=16mm of start,minimum width=3.0cm,fill=green!8] (connected)
    {CONNECTED\\gửi \codeword{register}};
  \node[block,left=of connected,minimum width=3.2cm] (run)
    {Mỗi 5 s gửi status\\liên tục drain UART};
  \node[decision,left=of run] (restart) {URL đổi / lỗi /\\Wi-Fi mất?};

  \draw[flow] (wait) -- (wifi);
  \draw[flow,loop above] (wifi) to node[above,font=\scriptsize]{không} (wifi);
  \draw[flow] (wifi) -- node[above,font=\scriptsize]{có} (resolve);
  \draw[flow] (resolve) -- (start);
  \draw[dataflow] (start) -- (connected);
  \draw[dataflow] (connected) -- (run);
  \draw[flow] (run) -- (restart);
  \draw[dataflow,bend left=15] (restart) to node[below,font=\scriptsize]{không} (run);
  \draw[warnflow,bend left=16] (restart) to node[above,font=\scriptsize]{có: destroy client} (wait);
\end{tikzpicture}}
\caption{Máy trạng thái vòng đời WebSocket}
\end{figure}

Client dùng timeout mạng và reconnect 15 giây, buffer 2048 B và keepalive.
\codeword{save_ws_url()} cập nhật URL và yêu cầu restart client, nhưng chỉ ghi
vào RAM. Sau reset, URL quay lại lựa chọn trong Kconfig/sdkconfig. Vì vậy thanh
cấu hình frontend hiện thay đổi được nguồn trong phiên chạy, chưa phải cấu hình
bền vững.

\subsection{Trình tự gửi dữ liệu đầu cuối}

\begin{figure}[p]
\centering
\resizebox{.98\textwidth}{!}{
\begin{tikzpicture}[x=2.75cm,y=.72cm,font=\scriptsize]
  \foreach \x/\name in {0/Node root,1/UART task,2/WS task,3/Backend,4/Dashboard}{
    \node[font=\bfseries] at (\x,0) {\name};
    \draw[gray,dashed] (\x,-.4) -- (\x,-12.3);
  }
  \draw[dataflow] (0,-1) -- node[above]{JSON kết thúc bằng LF} (1,-1);
  \draw[flow] (1,-2) -- node[above]{item tối đa 512 B} (2,-2);
  \draw[dataflow] (2,-3) -- node[above]{\codeword{uart_rx}} (3,-3);
  \draw[flow] (3,-4) -- node[above]{parse outer JSON} (3,-5);
  \draw[flow] (3,-5.5) -- node[above]{parse \codeword{payload}, lấy \codeword{M}, \codeword{p}} (3,-6.5);
  \draw[dataflow] (3,-7) -- node[above]{event dữ liệu node} (4,-7);
  \draw[dataflow] (2,-8.5) -- node[above]{\codeword{gateway_status}, mỗi 5 s} (3,-8.5);
  \draw[dataflow] (3,-9.5) -- node[above]{event trạng thái gateway} (4,-9.5);
  \draw[flow] (4,-11) -- node[above]{subscribe WS từ browser} (3,-11);
  \node[note,anchor=west] at (0,-12)
    {Hai loại bản tin phải được backend phân luồng theo trường \codeword{type}.};
\end{tikzpicture}}
\caption{Sequence diagram từ root đến dashboard}
\label{fig:sequence}
\end{figure}

\section{Nguồn tạo bản tin gateway\_status}

\begin{figure}[H]
\centering
\resizebox{\textwidth}{!}{
\begin{tikzpicture}[node distance=6mm and 9mm]
  \node[block,minimum width=3.0cm] (cpu) {FreeRTOS runtime stats\\CPU load};
  \node[block,below=of cpu,minimum width=3.0cm] (heap) {Internal heap\\RAM used/free};
  \node[block,below=of heap,minimum width=3.0cm] (adc) {ADC GPIO34\\pin voltage/\%};
  \node[block,below=of adc,minimum width=3.0cm] (timer) {\codeword{esp_timer}\\uptime};
  \node[block,below=of timer,minimum width=3.0cm] (net) {Wi-Fi driver/netif\\SSID, RSSI, IP, GW};

  \node[block,right=19mm of adc,minimum width=3.4cm,minimum height=5.4cm,fill=cyan!8] (mon)
    {\textbf{SystemMonitor}\\chu kỳ 1 s\\tạo snapshot\\\codeword{ui_metrics_t}\\
    ghi dưới mutex};
  \node[block,right=19mm of mon,minimum width=3.5cm,minimum height=3.1cm,fill=green!8] (json)
    {\textbf{WSHandle}\\chu kỳ 5 s\\đọc metrics + Wi-Fi\\tạo JSON
    \codeword{gateway_status}};

  \foreach \s in {cpu,heap,adc,timer,net} \draw[dataflow] (\s) -- (mon);
  \draw[dataflow] (mon) -- (json);
\end{tikzpicture}}
\caption{Nguồn dữ liệu tạo trạng thái gateway}
\label{fig:statussource}
\end{figure}

\begin{longtable}{p{4.4cm}p{4.2cm}p{5.3cm}}
\caption{Các trường của gateway\_status}\\
\toprule
\textbf{Trường} & \textbf{Nguồn} & \textbf{Ý nghĩa}\\
\midrule
\endfirsthead
\toprule
\textbf{Trường} & \textbf{Nguồn} & \textbf{Ý nghĩa}\\
\midrule
\endhead
\codeword{cpu_load_percent} & Runtime stats & Tải CPU trung bình, đơn vị \%.\\
\codeword{ram_used_kb}, \codeword{ram_used_percent} & Internal heap & RAM nội bộ đã dùng.\\
\codeword{battery_voltage_v}, \codeword{battery_percent} & ADC + PowerManager & Điện áp gói pin và mức pin ước lượng.\\
\codeword{power_source} & Điện áp pin & Heuristic: \codeword{unknown}, \codeword{battery}, \codeword{AC_or_USB}.\\
\codeword{chip_temp_c} & Cảm biến nội chip & Có thể là \codeword{null} nếu SoC không hỗ trợ.\\
\codeword{uptime_s} & esp\_timer & Số giây từ lần boot/reset.\\
\codeword{wifi_ssid}, \codeword{wifi_rssi} & Wi-Fi AP record & Mạng đang kết nối và RSSI dBm.\\
\codeword{sta_ip}, \codeword{sta_gateway} & STA netif & IPv4 của gateway và router.\\
\bottomrule
\end{longtable}

\section{Màn hình và dữ liệu cục bộ}

LVGL đọc snapshot metrics và cập nhật UI khoảng mỗi 250 ms. Các giá trị RX/TX,
IP, MAC rút gọn và mức sử dụng queue được cập nhật theo chu kỳ 1 giây.
SystemMonitor đo RTC, pin, SD, heap, CPU, FPS, uptime và nhiệt độ chip mỗi giây.

\begin{figure}[H]
\centering
\resizebox{.94\textwidth}{!}{
\begin{tikzpicture}[node distance=9mm and 12mm]
  \node[block,minimum width=3.2cm] (sources) {RTC / ADC / SD\\heap / runtime};
  \node[task,right=of sources,minimum width=3.4cm] (monitor) {system\_monitor\\1 s};
  \node[store,right=of monitor,minimum width=3.4cm] (metrics) {metrics snapshot\\mutex};
  \node[task,right=of metrics,minimum width=3.4cm] (lvgl) {lvgl\_task\\UI 250 ms};
  \node[block,below=9mm of lvgl,minimum width=3.4cm] (display) {Màn hình\\status + chart};
  \node[task,below=9mm of metrics,minimum width=3.4cm] (ws) {ws\_hdl\\status 5 s};
  \draw[dataflow] (sources) -- (monitor);
  \draw[dataflow] (monitor) -- (metrics);
  \draw[dataflow] (metrics) -- (lvgl);
  \draw[dataflow] (lvgl) -- (display);
  \draw[dataflow] (metrics) -- (ws);
\end{tikzpicture}}
\caption{Một snapshot metrics phục vụ cả màn hình và WebSocket}
\end{figure}

SPIFFS hiện lưu tài nguyên captive portal. Thẻ SD được khởi tạo và đo dung lượng,
nhưng đường nhận UART--WebSocket hiện không ghi telemetry vào SD, SQLite hay
MongoDB. Việc lưu dữ liệu dài hạn thuộc backend; gateway chỉ chuyển tiếp trong
RAM.

\section{Phân vùng flash và OTA}

\begin{figure}[H]
\centering
\resizebox{.94\textwidth}{!}{
\begin{tikzpicture}[x=1cm,y=1cm,font=\small]
  \draw[fill=orange!18] (0,0) rectangle (2.1,1.35);
  \draw[fill=yellow!25] (2.1,0) rectangle (4.2,1.35);
  \draw[fill=blue!18] (4.2,0) rectangle (8.2,1.35);
  \draw[fill=cyan!18] (8.2,0) rectangle (12.2,1.35);
  \draw[fill=green!18] (12.2,0) rectangle (14.5,1.35);
  \node[align=center] at (1.05,.68) {\textbf{NVS}\\\codeword{0x5000}};
  \node[align=center] at (3.15,.68) {\textbf{otadata}\\\codeword{0x2000}};
  \node[align=center] at (6.2,.68) {\textbf{ota\_0}\\\codeword{0x1C0000}};
  \node[align=center] at (10.2,.68) {\textbf{ota\_1}\\\codeword{0x1C0000}};
  \node[align=center] at (13.35,.68) {\textbf{SPIFFS}\\\codeword{0x20000}};
  \node[note,anchor=north] at (7.25,-.25)
    {Sơ đồ ưu tiên khả năng đọc; chiều rộng NVS, otadata và SPIFFS đã được phóng đại.};
\end{tikzpicture}}
\caption{Các partition flash và kích thước khai báo}
\end{figure}

\begin{figure}[H]
\centering
\resizebox{.94\textwidth}{!}{
\begin{tikzpicture}[node distance=8mm and 10mm]
  \node[block,minimum width=3.1cm] (server) {Server gửi\\version/lệnh OTA};
  \node[block,right=of server,minimum width=3.2cm] (handler) {WS handler\\chỉ đọc type/version};
  \node[block,right=of handler,minimum width=3.2cm,fill=red!7] (gap) {Điểm nối còn thiếu\\không gọi OTA};
  \node[block,right=of gap,minimum width=3.2cm] (fota) {FOTAManager\\mã phần lớn bị comment};
  \node[store,below=10mm of fota,minimum width=3.2cm] (part) {ota\_0 / ota\_1};
  \draw[flow] (server) -- (handler);
  \draw[warnflow] (handler) -- (gap);
  \draw[warnflow] (gap) -- (fota);
  \draw[warnflow] (fota) -- (part);
\end{tikzpicture}}
\caption{Trạng thái triển khai OTA: partition có sẵn nhưng pipeline chưa nối}
\label{fig:ota}
\end{figure}

Không nên hiểu sự tồn tại của hai partition OTA là FOTA đã hoạt động. Trong luồng
hiện tại, WebSocket chỉ lưu/log \codeword{type} và \codeword{version}; không gọi
hàm tải firmware. Khi hoàn thiện cần thêm xác thực manifest, kiểm tra phiên bản,
chữ ký/hash, chính sách certificate và rollback.

\section{Tài nguyên FreeRTOS}

\begin{table}[ht]
\centering
\caption{Các task chính nhìn thấy trong mã}
\begin{tabularx}{\textwidth}{p{3.2cm}p{2cm}p{1.5cm}X}
\toprule
\textbf{Task} & \textbf{Stack} & \textbf{Priority} & \textbf{Vai trò}\\
\midrule
\codeword{system_monitor} & 3072 B & 5 & Metrics, pin, RTC, SD, heap, CPU; pin core 0.\\
\codeword{lvgl_task} & 4096 B & 4 & Render và cập nhật màn hình; pin core 1.\\
\codeword{ws_hdl} & 4096 B & 5 & Quản lý client và gửi dữ liệu. Client WS nội bộ yêu cầu stack 6144 cho task của thư viện.\\
\codeword{uart_to_node} & 4096 B mặc định & 5 & Bắt tay và đọc/ghi UART.\\
\codeword{wifi_connect} & 4096 B & 5 & Áp dụng credential mới.\\
\codeword{wifi_mgr} & 4096 B & 5 & Giám sát STA/AP và reconnect.\\
\codeword{dns_server} & 2560 B & 5 & DNS captive portal khi AP hoạt động.\\
\bottomrule
\end{tabularx}
\end{table}

SystemMonitor cấp phát mảng \codeword{TaskStatus_t} bằng \codeword{malloc()} mỗi
giây để tính CPU load. Có thể giảm phân mảnh heap bằng buffer tái sử dụng hoặc
chu kỳ đo dài hơn.

\section{Rủi ro và đề xuất bảo trì}

\begin{enumerate}
  \item \textbf{Bí mật trong mã nguồn:} auth key ESP Insights và một số mật khẩu
  Wi-Fi đang viết trực tiếp trong tệp C. Cần thu hồi/đổi khóa đã lộ và chuyển sang
  provisioning, biến môi trường build hoặc kho bí mật.
  \item \textbf{URL WebSocket không bền vững:} thay đổi qua web chỉ lưu RAM. Nên
  lưu URL hợp lệ vào NVS và tải trước khi tạo client.
  \item \textbf{Không có lưu trữ offline:} mất WS đồng nghĩa bỏ telemetry UART.
  Có thể bổ sung ring buffer trong RAM hoặc journal trên SD, kèm giới hạn dung lượng.
  \item \textbf{Tách JSON hai lớp:} backend phải parse
  \codeword{uart_rx.payload}. Nên định nghĩa schema và test contract cho trường
  \codeword{M}, \codeword{p}, \codeword{seq} và \codeword{err}.
  \item \textbf{Danh sách Wi-Fi boot làm chậm khởi động:} nhiều lần chờ 20 giây
  có thể kéo dài boot. Nên scan trước, chỉ thử SSID nhìn thấy và giới hạn ngân sách thời gian.
  \item \textbf{FOTA chưa hoàn chỉnh:} không mở chức năng OTA cho sản phẩm trước
  khi có xác thực, rollback và kiểm thử mất điện.
  \item \textbf{mDNS phụ thuộc mạng LAN:} \codeword{systemmsems.local} chỉ phân
  giải khi máy backend quảng bá mDNS và router không chặn multicast. Nên cho phép
  nhập IP/hostname thay thế và hiển thị lỗi phân giải rõ ràng.
\end{enumerate}

\section{Quy trình build và kiểm thử}

\begin{lstlisting}[style=textstyle,caption={Build và flash firmware}]
cd MeshGateWay
idf.py set-target esp32
idf.py menuconfig
idf.py build
idf.py -p COMx flash monitor
\end{lstlisting}

Các ca kiểm thử tối thiểu:
\begin{itemize}
  \item boot khi có/không có Wi-Fi; kiểm tra AP, DNS captive portal và
  \codeword{meshgateway.local};
  \item local \codeword{ws://systemmsems.local:9090} và server
  \codeword{wss://...}; xác minh log hiển thị \codeword{plain} hoặc \codeword{TLS};
  \item ngắt UART quá 2 giây để kiểm tra quay lại \codeword{START_ROOT};
  \item gửi JSON dài, chia qua nhiều chunk và xác minh backend nhận đúng một dòng;
  \item làm đầy queue hoặc ngắt WS để đo số bản tin bị mất;
  \item đối chiếu \codeword{gateway_status} với RSSI, IP, pin và heap thực tế;
  \item reset sau khi đổi URL để xác nhận giới hạn lưu RAM hiện tại.
\end{itemize}

\section{Kết luận}

MeshGateWay được tổ chức thành ba đường xử lý chính: UART nhận dữ liệu từ root,
SystemMonitor tạo metrics nội tại, và WebSocket hợp nhất hai nguồn để gửi backend.
FreeRTOS queue tách task UART khỏi task mạng; mutex tạo snapshot ổn định cho UI
và status. Local dùng \codeword{ws://} qua mDNS, còn server dùng
\codeword{wss://} với TLS. Các giới hạn quan trọng nhất của phiên bản hiện tại là
không lưu telemetry offline, URL WebSocket chỉ lưu RAM, FOTA chưa nối và còn bí
mật được viết trực tiếp trong mã nguồn.

\appendix
\section{Thứ tự đọc mã đề xuất}

\begin{longtable}{p{6.1cm}p{7.9cm}}
\toprule
\textbf{Tệp} & \textbf{Nội dung nên đọc}\\
\midrule
\endfirsthead
\toprule
\textbf{Tệp} & \textbf{Nội dung nên đọc}\\
\midrule
\endhead
\pathcode{main/MeshGateWay.c} & Trình tự khởi động và quan hệ tổng thể.\\
\pathcode{components/core/Datamanager/Datamanager.h} & Toàn bộ kiểu trạng thái và kích thước queue.\\
\pathcode{components/core/UartToNode/UartToNode.c} & Máy trạng thái bắt tay và uplink UART.\\
\pathcode{components/core/WebSocKetHandle/WSHandle.c} & URL, mDNS, WS/WSS và schema bản tin.\\
\pathcode{components/core/WifiManager/wifi_manager_core.c} & Event Wi-Fi, danh sách boot và mDNS.\\
\pathcode{components/core/WifiManager/wifi_manager_tasks.c} & Reconnect và chuyển STA/AP.\\
\pathcode{components/core/WifiManager/wifi_manager_http.c} & Endpoint cấu hình web.\\
\pathcode{components/core/SystemMonitor/SystemMonitor.c} & Cách tạo metrics.\\
\pathcode{components/core/ScreenManager/ScreenManager.c} & Cách metrics được đưa lên LVGL.\\
\pathcode{partitions.csv} & Phân vùng NVS, hai OTA slot và SPIFFS.\\
\bottomrule
\end{longtable}

\end{document}
